\documentclass[10pt,a4paper]{amsart}
\usepackage[margin=19mm]{geometry}
\usepackage{amsmath,amssymb,mathtools}
\usepackage[scr=rsfs,cal=euler]{mathalfa}
\usepackage{xfrac}
\usepackage[unicode,hidelinks,bookmarks=false]{hyperref}
\newcommand{\ie}{i.e.\ }
\newcommand{\cf}{cf.\ }
\newcommand{\resp}{respectively\ }
\newcommand{\Subsection}{\subsection}
\providecommand{\nobreakdash}{\nobreak\hbox{-}}
\setlength{\emergencystretch}{2em}
\multlinegap0pt
\usepackage{bm}
\usepackage{bbm}
\usepackage{dsfont}
\usepackage{xspace}
\usepackage{cancel}
\usepackage{mathtools}
\usepackage{amsthm}
\makeatletter
\@ifundefined{theorem}{\newtheorem{theorem}{Theorem}}{}
\@ifundefined{corollary}{\newtheorem{corollary}[theorem]{Corollary}}{}
\@ifundefined{proposition}{\newtheorem{proposition}[theorem]{Proposition}}{}
\makeatother
\newcommand{\F}{\mathbb{F}}
\title[Shift-invariant transformations and almost liftings]{Shift-invariant transformations and almost liftings}
\author{Jan Kristian Haugland, Tron Omland}
\date{Source: arXiv:2407.11931v3 (CC BY 4.0)}
\begin{document}
\maketitle

\noindent\emph{Source and licence.} Shift-invariant transformations and almost liftings. Version arXiv:2407.11931v3 (CC BY 4.0), distributed under the Creative Commons Attribution 4.0 International licence, as recorded in the arXiv licence metadata for this exact version and shown on the abstract page https://arxiv.org/abs/2407.11931v3. Full rights notice: licenses/arxiv-2407.11931-CC-BY-4.0.txt.\par\smallskip

\noindent\textit{Editorial scope.} Counted unit: the Corollary whose source label is \texttt{permutive-bound} in the article (source0.tex lines 251--260, source label \texttt{permutive-bound}), delivered together with the article's own complete original proof of it (source0.tex lines 262--293, one \texttt{proof} environment of 32 lines). No mathematics is written, repaired or re-derived: statement and proof are the article's own text line for line. Required context retained uncounted: permutive (lines 220--222). Every definition, display and label of those lines is retained unchanged. Omitted from this package and not redistributed: 1--219, 223--250, 261--261, 294--1043. Nothing inside a retained excerpt is omitted or altered: every retained body is delivered exactly as the article's lines are written, so no omission receipt of any kind accompanies this package. Citations: the retained excerpts carry no citation command at all, so no bibliography entry of the article is transcribed and no reference is redistributed. Reference adaptation: none was needed and none was made; every reference inside the delivered unit resolves inside it, so no unresolved reference or citation is produced. Printed slips kept literal: no printed word, symbol, hypothesis or number of the counted statement or of its proof was altered. Layout and class: the article's own class is \texttt{amsart} with packages including inputenc, fontenc, lmodern, babel, amsmath, amssymb, amsthm, geometry, microtype, hyperref, url, tikz-cd, colortbl, ulem; the wrapper is the lane's pinned \texttt{amsart} editorial wrapper at 10pt on A4, which restates the theorem-like environments the retained text needs and adds the article's own macro definitions that the retained text needs (\texttt{\textbackslash F}), with extra wrapper packages (bm, bbm, dsfont, xspace, cancel, mathtools). The delivered numbering is the unit's own. Assets: no retained span contains a figure environment, a graphics-inclusion command, a PDF-page inclusion, a table or a TikZ picture; no image file is copied into this package, the package declares no asset dependency and no image file is redistributed with it. Licence: version arXiv:2407.11931v3 of this article is distributed under the Creative Commons Attribution 4.0 International licence, as recorded in the arXiv licence metadata for this exact version and shown on the abstract page https://arxiv.org/abs/2407.11931v3. A full rights notice is delivered at licenses/arxiv-2407.11931-CC-BY-4.0.txt. Characters: every retained excerpt is pure ASCII, so the delivered engine, XeLaTeX with its default Latin Modern text font, reports no missing character.
	\begin{proposition}\label{permutive}
		For any two Boolean functions $h, h'\colon\F_2^{k-1} \to \F_2$, where $h(x)$ depends on $x_1$ and $h'(x)$ depends on $x_{k-1}$, the permutive functions $f, g\colon\F_2^k \to \F_2$ given by $f(x_1, \dotsc, x_k) = h(x_1, \dotsc, x_{k-1}) \oplus x_k$ and $g(x_1, \dotsc, x_k) = x_1 \oplus h'(x_2, \dotsc, x_k)$ are potential $(k, n)$-liftings for all $n \geq k$.
	\end{proposition}

    \begin{corollary}\label{permutive-bound}
    The number of elementary equivalence classes of permutive Boolean functions of diameter $k$ is equal to
    \[
    \frac{1}{8}
    \begin{cases}
        2 \cdot 2^{2^{k-1}} + 2^{2^{k-2}} + 2 \cdot 2^{2^{k-3} + 2^{\frac{k}{2}-2}} - 6 \cdot 2^{\lfloor 2^{k-3} \rfloor} &\quad\text{if }k \equiv 0\, (\operatorname{mod} 2)\\
        2 \cdot 2^{2^{k-1}} + 2^{2^{k-2}} + 2^{2^{k-3} + 2^{\frac{k-3}{2}}} - 4 \cdot 2^{2^{k-3}} &\quad\text{if }k \equiv 1\, (\operatorname{mod} 2)
    \end{cases}
    \]
    \end{corollary}

    \begin{proof}
    Let $T_k$ denote the set of potential $(k, n)$-liftings given by Lemma~\ref{permutive}. Let $T^r_k = \{ f \in T_k : f(x) = f \circ r(x) \oplus \kappa\}$ for some $\kappa$, and similarly $T^c_k = \{ f \in T_k : f(x) = f \circ c(x) \oplus \kappa\}$ and $T^{r \circ c}_k = \{ f \in T_k : f(x) = f \circ r \circ c(x) \oplus \kappa\}$ for some $\kappa$. The expression $\lvert T_k\rvert + \lvert T^r_k\rvert + \lvert T^c_k\rvert + \lvert T^{r \circ c}_k\rvert$ counts each function with 8 elements in its equivalence class once, each function with 4 elements in its equivalence class twice, and each function with 2 elements in its equivalence class four times. Thus, the number of equivalence classes of $T_k$ is $$\frac{\lvert T_k\rvert + \lvert T^r_k\rvert + \lvert T^c_k\rvert + \lvert T^{r \circ c}_k\rvert}{8}.$$
    
    The number of functions of the form $h(x_1, \dotsc, x_{k-1}) \oplus x_k$ that depend on $x_1$ is $2^{2^{k-1}} - 2^{2^{k-2}}$, and similarly for functions of the form $x_1 \oplus h'(x_2, \dotsc, x_k)$ that depend on $x_k$. There are $2^{2^{k-2}}$ functions in the intersection, i.e., of the form $x_1 \oplus h(x_2, \dotsc,  x_{k-1}) \oplus x_k$. Thus, $\lvert T_k\rvert = 2 \cdot 2^{2^{k-1}} - 3 \cdot 2^{2^{k-2}}$.
    
    The functions in $T^r_k$ must have $\kappa=0$ since $r$ fixes some inputs, and have the form $f(x) = x_1 \oplus h(x_2, \dotsc, x_{k-1}) \oplus x_k$. The expression $2^{k-2} + 2^{\lfloor \frac{k-1}{2}\rfloor}$ counts the $x$'s for which $h(x) = h \circ r(x)$ twice and those for which $h(x) \neq h \circ r(x)$ once. Thus,
    \[
    \lvert T^r_k\rvert=
    \begin{cases}
        2^{\frac{2^{k-2} + 2^{\frac{k}{2}-1}}{2}} &\quad\text{if }k \equiv 0\, (\operatorname{mod} 2)\\
        2^{\frac{2^{k-2} + 2^{\frac{k-1}{2}}}{2}} &\quad\text{if }k \equiv 1\, (\operatorname{mod} 2)
    \end{cases}
    \]
    The functions in $T^c_k$ could have $\kappa$ equal to either 0 or 1, except for $k=2$, in which case we must have $f(x) = x_1 \oplus x_2 \oplus $ some constant which implies $\kappa=0$. Like in the derivation of $\lvert T_k\rvert$ above, we have $$\lvert T^c_k\rvert = 2 \left( 2 \cdot 2^{2^{k-2}} - 3 \cdot 2^{\lfloor 2^{k-3} \rfloor} \right)$$ (which happens to give the correct number also for $k=2$).
    
    The functions in $T^{r \circ c}_k$ must have $\kappa=0$ if $k$ is even, otherwise it can be either 0 or 1, and they must have the form $f(x) = x_1 \oplus h(x_2, \dotsc, x_{k-1}) \oplus x_k$. Thus,
    \[
    \lvert T^{r \circ c}_k\rvert=
    \begin{cases}
        2^{\frac{2^{k-2}+2^{\frac{k}{2}-1}}{2}} &\quad\text{if }k \equiv 0\, (\operatorname{mod} 2)\\
        2 \cdot 2^{\frac{2^{k-2}}{2}} &\quad\text{if }k \equiv 1\, (\operatorname{mod} 2)
    \end{cases}
    \]
    Thus, the number of classes is
    \[
    \frac{1}{8}
    \begin{cases}
        2 \cdot 2^{2^{k-1}} + 2^{2^{k-2}} + 2 \cdot 2^{2^{k-3} + 2^{\frac{k}{2}-2}} - 6 \cdot 2^{\lfloor 2^{k-3} \rfloor} &\quad\text{if }k \equiv 0\, (\operatorname{mod} 2)\\
        2 \cdot 2^{2^{k-1}} + 2^{2^{k-2}} + 2^{2^{k-3} + 2^{\frac{k-3}{2}}} - 4 \cdot 2^{2^{k-3}} &\quad\text{if }k \equiv 1\, (\operatorname{mod} 2)
    \end{cases}
    \]
    \end{proof}

\end{document}
